\clearpage\thispagestyle{empty}
\pdfbookmark[0]{Sources for Gauss’s Observations of Iris (September 1847)}{iris1847-sources}
\begin{center}\textbf{Sources for Gauss’s Observations of Iris (September 1847)}\end{center}
\noindent This AI translation follows the report by Carl Friedrich Gauss on p.~473 of \href{https://gdz.sub.uni-goettingen.de/id/PPN236007467}{\textit{Werke}, vol.~VI} (Göttingen, 1874), dated 7 September 1847. The first printing appeared in \textit{Nachrichten von der Georg-Augusts-Universität und der Königl. Gesellschaft der Wissenschaften zu Göttingen}, no.~11 of \href{https://images.p.sub.uni-goettingen.de/iiif/image/gdz:PPN668664029_0003:00000163/full/1800,/0/default.jpg}{13 September 1847}, \href{https://images.p.sub.uni-goettingen.de/iiif/image/gdz:PPN668664029_0003:00000168/full/1800,/0/default.jpg}{pp.~166}, \href{https://images.p.sub.uni-goettingen.de/iiif/image/gdz:PPN668664029_0003:00000169/full/1800,/0/default.jpg}{167} and \href{https://images.p.sub.uni-goettingen.de/iiif/image/gdz:PPN668664029_0003:00000170/full/1800,/0/default.jpg}{168}. The report begins under the heading of the Royal Society of Sciences on p.~166.
\par\smallskip\noindent John Russell Hind discovered Iris. Gauss reports his own observations and distinguishes Goldschmidt’s elements, based on observations by Gauss and Challis, from Gould’s separately calculated orbit. The table notation and punctuation here follow the collected edition. Translation and typesetting produced by AI.
\clearpage
